\documentclass[11pt]{article}
\usepackage[margin=1in]{geometry}
\usepackage{amsmath,amssymb,booktabs}
\usepackage{hyperref}
\usepackage{enumitem}
\usepackage[T1]{fontenc}
\usepackage[utf8]{inputenc}

\title{Indirect calorimetry to gene-set extractors for CalR-compatible datasets}
\author{}
\date{March 2026}

\begin{document}
\maketitle

\begin{abstract}
Indirect calorimetry is a whole-animal physiology assay, not a direct molecular assay. A CalR-uploadable dataset therefore reaches genes only through phenotype routing or reference-profile retrieval. This note describes the two current extractors in \texttt{geneset-extractors}: \texttt{calr\_ontology\_mapper} and \texttt{calr\_profile\_query}, together with the CalR processing rules they preserve.
\end{abstract}

\section{Mental model}
Let $y_{ivt}$ denote the measured value of variable $v$ for animal $i$ at time bin $t$. The extractor does not map raw bins directly to genes. It first computes selected-window summary features, then estimates contrast-aware physiology signatures, then routes those signatures to genes.

The core pipeline is:
\begin{enumerate}[leftmargin=2em]
\item choose the active analysis window and photoperiod,
\item summarize each animal into feature vector $s_i$,
\item estimate a contrast or grouped signature $q$ using CalR-style covariate logic,
\item either map $q$ to phenotype terms and then genes, or compare $q$ to a reference library and route retrieved evidence to genes,
\item select and normalize the resulting gene program.
\end{enumerate}

\section{CalR compatibility rules}
The extractor preserves the following constraints from CalR-style analysis:
\begin{itemize}[leftmargin=2em]
\item the selected analysis window is not silently the full trace,
\item session metadata is preferred and recorded as \texttt{session\_mode=explicit},
\item missing session metadata is allowed only in exploratory mode with strong warnings,
\item mass-dependent variables are not naively divided by body weight,
\item \texttt{ee} and \texttt{rer} are used as reported when present,
\item interaction with the selected mass covariate is tested before falling back to a simpler model,
\item exclusions are explicit and auditable.
\end{itemize}

Practical covariate priority is:
\begin{enumerate}[leftmargin=2em]
\item explicit user covariate,
\item \texttt{Lean.Mass},
\item \texttt{Total.Mass},
\item \texttt{subject.mass},
\item no covariate only as an exploratory fallback.
\end{enumerate}

\section{Selected-window feature derivation}
For each animal, define the selected bin set
\begin{equation}
\mathcal{W}_i = \{ t : m_i(t) = 1 \},
\end{equation}
where $m_i(t)$ reflects the selected window and any explicit exclusions.

Whole-window averages are
\begin{equation}
\bar{y}^{(W)}_{iv} = \frac{1}{|\mathcal{W}_i|} \sum_{t \in \mathcal{W}_i} y_{ivt}.
\end{equation}

Light and dark summaries are computed on the selected-window subsets:
\begin{equation}
\bar{y}^{(L)}_{iv} = \frac{1}{|\mathcal{W}^{(L)}_i|} \sum_{t \in \mathcal{W}^{(L)}_i} y_{ivt},
\qquad
\bar{y}^{(D)}_{iv} = \frac{1}{|\mathcal{W}^{(D)}_i|} \sum_{t \in \mathcal{W}^{(D)}_i} y_{ivt}.
\end{equation}

Day-night contrast:
\begin{equation}
\Delta^{(D-L)}_{iv} = \bar{y}^{(D)}_{iv} - \bar{y}^{(L)}_{iv}.
\end{equation}

For cumulative variables:
\begin{equation}
A_{iv} = \sum_{t \in \mathcal{W}_i} y_{ivt}\Delta t.
\end{equation}

For circadian structure, the first harmonic is summarized by
\begin{equation}
a_{iv} = \frac{2}{|\mathcal{W}_i|}\sum_{t \in \mathcal{W}_i} y_{ivt}\cos(2\pi t / 24),
\qquad
b_{iv} = \frac{2}{|\mathcal{W}_i|}\sum_{t \in \mathcal{W}_i} y_{ivt}\sin(2\pi t / 24),
\end{equation}
with amplitude
\begin{equation}
r_{iv} = \sqrt{a_{iv}^2 + b_{iv}^2}.
\end{equation}

\section{Variable-class-aware contrast estimation}
Mass-independent variables are treated without the mass covariate:
\begin{equation}
s_{ik} = \alpha_k + z_i^\top \gamma_k + \varepsilon_{ik}.
\end{equation}

Mass-dependent variables use a CalR-style ANCOVA/GLM split:
\begin{equation}
s_{ik} = \alpha_k + x_i^\top \beta_k + z_i^\top \gamma_k + (x_i \otimes z_i)^\top \delta_k + \varepsilon_{ik},
\end{equation}
where $x_i$ is the selected mass/body-composition covariate and $z_i$ encodes the experimental contrast.

Implementation rule:
\begin{enumerate}[leftmargin=2em]
\item test the interaction term,
\item retain it when non-negligible,
\item otherwise drop it and use the simpler ANCOVA-style effect.
\end{enumerate}

The resulting standardized feature effect is treated as the query signature component $q_k$.

\section{Program decomposition}
The extractor does not force one omnibus physiology output. Instead it emits program-specific views:
\begin{itemize}[leftmargin=2em]
\item \texttt{global}
\item \texttt{thermogenesis}
\item \texttt{substrate\_use}
\item \texttt{intake\_balance}
\item \texttt{activity\_circadian}
\end{itemize}

Each program masks the feature space before routing.

\section{Extractor 1: \texttt{calr\_ontology\_mapper}}
Let $p$ index phenotype terms with signed feature templates $T_{pk}$. The raw term score is
\begin{equation}
z_p(q) = \sum_{k=1}^{K} T_{pk} q_k.
\end{equation}

A nonnegative normalized term score is
\begin{equation}
u_p(q) = \frac{\max(0, z_p(q))}{\|T_p\|_1 + \epsilon}.
\end{equation}

Let $H_{pg}$ be a phenotype-to-gene adjacency from phenotype-gene resources. Then the routed ontology score is
\begin{equation}
S^{\mathrm{onto}}_g(q) = \sum_p \omega_p u_p(q) H_{pg},
\end{equation}
where $\omega_p$ downweights generic or high-degree terms.

The current implementation emits:
\begin{itemize}[leftmargin=2em]
\item \texttt{mode=core}: direct routed phenotype support
\item \texttt{mode=expanded}: small hierarchy-aware term expansion when a hierarchy is available
\end{itemize}

\section{Extractor 2: \texttt{calr\_profile\_query}}
Let $x_r$ be a reference feature profile aligned to the same feature schema as the query. After optional centering/scaling, similarity is
\begin{equation}
\mathrm{sim}(q,r) = \frac{\tilde{q}^\top \tilde{x}_r}{\|\tilde{q}\|_2 \|\tilde{x}_r\|_2}
\end{equation}
for cosine similarity, or Pearson correlation when that metric is selected.

The current weighted reference evidence is
\begin{equation}
w_r(q) = a_r \, h_r \, \max(0, \mathrm{sim}(q,r) - \tau)^\gamma \, \pi_r,
\end{equation}
where:
\begin{itemize}[leftmargin=2em]
\item $a_r$ is any implicit reference quality term,
\item $h_r$ is the chosen hubness penalty,
\item $\tau$ is the similarity floor,
\item $\gamma$ is the similarity power,
\item $\pi_r$ is a provenance penalty for mismatched query/reference metadata.
\end{itemize}

With one primary perturbed gene per reference, the direct profile-query score is
\begin{equation}
S^{\mathrm{prof}}_g(q) = \sum_r w_r(q)\mathbf{1}\{g_r = g\}.
\end{equation}

The current implementation is intentionally conservative: it emits \texttt{mode=core} only and warns when retrieval confidence is low or when query/reference provenance looks mismatched.

\section{Reference bundles}
`workflows calr_prepare_reference_bundle` packages:
\begin{itemize}[leftmargin=2em]
\item reference profiles,
\item reference metadata,
\item feature schema,
\item feature stats,
\item ontology resources,
\item bundle manifest and summary,
\item optional tarball distribution artifact.
\end{itemize}

Packaged ontology defaults are mouse-first. Human profile-query use is supported when explicit human reference resources are supplied.

\section{Limitations}
Current v1 scope is:
\begin{itemize}[leftmargin=2em]
\item straightforward grouped or chronic designs,
\item CalR-standardized inputs,
\item phenotype routing and reference-profile retrieval,
\item no dedicated challenge-response or crossover modeling yet.
\end{itemize}

The extractor is therefore best described as a CalR-compatible phenotype-to-gene family, not a general-purpose parser for every vendor-native calorimetry export.

\end{document}
